% =========================================================
% 18. sorted() ORM Method
% =========================================================
\section{\texttt{sorted()} ORM Method}

\begin{frame}[standout]
  \texttt{sorted()} ORM Method
\end{frame}

\begin{frame}[fragile]{\texttt{SORTED()} METHOD}
  \begin{itemize}
    \item The \texttt{sorted()} method returns a recordset ordered in Python.
    \item It sorts an existing recordset; it does not query the database again.
    \item The method returns a new sorted recordset.
  \end{itemize}
  \begin{block}{Method Syntax}
\begin{lstlisting}
def sorted(self, key=None, reverse=False):
    return super().sorted(key=key, reverse=reverse)
\end{lstlisting}
    Important parts: \texttt{key}, \texttt{reverse}.
  \end{block}
\end{frame}

\begin{frame}[fragile]{BREAKING DOWN THE METHOD SIGNATURE}
  \begin{itemize}
    \item \texttt{self}: The recordset to sort.
    \item \texttt{key}: A field name string, or a function that returns the sort value.
    \item \texttt{reverse}: \texttt{False} for ascending, \texttt{True} for descending.
  \end{itemize}
\begin{lstlisting}
students.sorted('name')
students.sorted(key=lambda s: s.age, reverse=True)
\end{lstlisting}
\end{frame}

\begin{frame}[fragile]{WHAT DOES \texttt{SORTED()} RETURN?}
  \begin{itemize}
    \item It returns a recordset in the new order.
  \end{itemize}
\begin{lstlisting}
students = self.env['student.student'].browse([17, 15, 16])
ordered = students.sorted('name')
print(ordered)
\end{lstlisting}
  Output is (example):
\begin{lstlisting}
student.student(15, 16, 17)
\end{lstlisting}
\end{frame}

\begin{frame}[fragile]{SORTING BY FIELD OR LAMBDA}
\begin{lstlisting}
by_name = students.sorted('name')
by_age_desc = students.sorted('age', reverse=True)
by_email = students.sorted(key=lambda s: s.email or '')
by_class_then_name = students.sorted(
    key=lambda s: (s.classroom_id.name or '', s.name or '')
)
\end{lstlisting}
  \begin{itemize}
    \item Use a lambda when you need multi-key or computed sorting.
  \end{itemize}
\end{frame}

\begin{frame}[fragile]{\texttt{SORTED()} VS \texttt{SEARCH(... ORDER=)}}
  \begin{itemize}
    \item \texttt{search(..., order='age desc')} sorts in SQL while fetching.
    \item \texttt{sorted()} sorts an already loaded recordset in Python.
  \end{itemize}
\begin{lstlisting}
# SQL ordering
students = self.env['student.student'].search([], order='age desc')

# Python ordering
students = students.sorted('age', reverse=True)
\end{lstlisting}
\end{frame}

\begin{frame}[fragile]{METHOD CALL}
  \textbf{Method Call}
\begin{lstlisting}
ordered_students = students.sorted(key='age', reverse=True)
\end{lstlisting}
  \begin{itemize}
    \item Here, \texttt{.sorted(...)} is the method call.
    \item Return type reminder: \textbf{recordset}.
  \end{itemize}
\end{frame}
